% Options for packages loaded elsewhere
\PassOptionsToPackage{unicode}{hyperref}
\PassOptionsToPackage{hyphens}{url}
%
\documentclass[
  10pt,
]{article}
\usepackage{amsmath,amssymb}
\usepackage{iftex}
\ifPDFTeX
  \usepackage[T1]{fontenc}
  \usepackage[utf8]{inputenc}
  \usepackage{textcomp} % provide euro and other symbols
\else % if luatex or xetex
  \usepackage{unicode-math} % this also loads fontspec
  \defaultfontfeatures{Scale=MatchLowercase}
  \defaultfontfeatures[\rmfamily]{Ligatures=TeX,Scale=1}
\fi
\usepackage{lmodern}
\ifPDFTeX\else
  % xetex/luatex font selection
\fi
% Use upquote if available, for straight quotes in verbatim environments
\IfFileExists{upquote.sty}{\usepackage{upquote}}{}
\IfFileExists{microtype.sty}{% use microtype if available
  \usepackage[]{microtype}
  \UseMicrotypeSet[protrusion]{basicmath} % disable protrusion for tt fonts
}{}
\makeatletter
\@ifundefined{KOMAClassName}{% if non-KOMA class
  \IfFileExists{parskip.sty}{%
    \usepackage{parskip}
  }{% else
    \setlength{\parindent}{0pt}
    \setlength{\parskip}{6pt plus 2pt minus 1pt}}
}{% if KOMA class
  \KOMAoptions{parskip=half}}
\makeatother
\usepackage{xcolor}
\usepackage[margin=0.85in]{geometry}
\usepackage{longtable,booktabs,array}
\usepackage{calc} % for calculating minipage widths
% Correct order of tables after \paragraph or \subparagraph
\usepackage{etoolbox}
\makeatletter
\patchcmd\longtable{\par}{\if@noskipsec\mbox{}\fi\par}{}{}
\makeatother
% Allow footnotes in longtable head/foot
\IfFileExists{footnotehyper.sty}{\usepackage{footnotehyper}}{\usepackage{footnote}}
\makesavenoteenv{longtable}
\setlength{\emergencystretch}{3em} % prevent overfull lines
\providecommand{\tightlist}{%
  \setlength{\itemsep}{0pt}\setlength{\parskip}{0pt}}
\setcounter{secnumdepth}{-\maxdimen} % remove section numbering
\usepackage{amsmath,amssymb,amsthm}
\usepackage{longtable}
\ifLuaTeX
  \usepackage{selnolig}  % disable illegal ligatures
\fi
\usepackage{bookmark}
\IfFileExists{xurl.sty}{\usepackage{xurl}}{} % add URL line breaks if available
\urlstyle{same}
\hypersetup{
  pdftitle={P16.5: Prescribed divisorial references in all big rational classes},
  pdfauthor={Fresh proof reconstruction relative to the exact ample grant G0},
  hidelinks,
  pdfcreator={LaTeX via pandoc}}

\title{P16.5: Prescribed divisorial references in all big rational
classes}
\author{Fresh proof reconstruction relative to the exact ample grant G0}
\date{5 October 2026}

\begin{document}
\maketitle

\section{Abstract}\label{abstract}

Over every algebraically closed trivially valued field of characteristic
zero, let \(X\) be a smooth integral projective variety and let \(L\) be
any big rational line bundle. Fix, before a target measure is chosen, an
actual decomposition \(L=A+D\), where \(A\) is an ample rational line
bundle and \(D\) is an effective rational Cartier divisor on \(X\). We
reconstruct the entire genuine cone dominated by the canonical section
reference \(\sigma_D\): it is exactly
\(\sigma_D+\operatorname{PSH}(X,A)\), with a unique actual ample
quotient, including its values on the reference divisor. The
reconstruction uses one defining FS net, rational rounding, a finite
polytope calculation, elementary convexity on scaling rays, literal
section division, and compactness exclusively in a fixed actual ample
cone. It does not use the false general extended-valued flattening lemma
in BJ v4.

The actual local metric isometry off \(D\), followed by reference
floors, constructs its genuine nonpluripolar mixed intersection operator
on the original \(X\). The raw mass is \((A^n)\), explicitly tied to the
prescribed reference. Relative energy, the intrinsic target class,
two-sided outer derivatives in every continuous original direction,
comparison, and an attained equation with uniqueness modulo constants
follow with their precise original-carrier interpretations. The
separately proved zero-volume range is the actual trivial bundle, whose
genuine cone consists of constants, whose operator and target mass are
zero, and whose energy is zero.

This is an unverified proof draft relative to the task's exact
actual-ample theorem G0. No result is inferred from a historical
verdict, hash, compilation, or finite test. The theorem covers every big
rational \(L\) with each fixed datum above. It does not construct the
minimal-singularity equation of mass \(\operatorname{vol}(L)\),
arbitrary reference types, or the equation theory for every nonbig
pseudoeffective class. These residual interpretations of original P16.5
remain explicit in the final scope assessment.

\section{1. Exact domain, reference and
statement}\label{exact-domain-reference-and-statement}

Let \(k\) be algebraically closed, trivially valued, and of
characteristic zero, with no restriction on cardinality. Let \(X\) be
smooth, integral, projective over \(k\), of dimension \(n\geq1\). All
analytic functions below are actual functions on the compact Hausdorff
Berkovich space \(X^{\rm an}\). Rational bundles mean elements of
\(\operatorname{Pic}(X)\otimes\mathbb Q\), cleared by positive integral
tensor powers whenever sections occur. Potentials are measured from the
canonical trivial metric on the indicated actual bundle. All measures
are positive Radon measures unless called signed. Intersections use
\([X]\) and have no probability factor.

For an actual rational bundle \(B\), a generically finite real FS
function is \[f=m^{-1}\max_j\{\log\|s_j\|_{\rm triv}+\lambda_j\},\qquad
s_j\in H^0(X,mB)\setminus\{0\},\quad\lambda_j\in\mathbb R.\tag{1.1}\]
The displayed family need not generate \(mB\). It generates exactly when
\(f\) is finite everywhere. Nonzero sections suffice for generic
finiteness on integral \(X\). The canonical norms lie in \([0,1]\), and
section logarithms take the value \(-\infty\) on their zero schemes. The
trivial valuation of \(k(X)\) is denoted \(v_0\).

For \(\theta\in N^1(X)\) we use the literal defining cone:
\(\phi\in\operatorname{PSH}_{\rm def}(\theta)\) means that \(\phi\) is
generically finite and is the all-point limit of one decreasing net of
functions (1.1) on actual rational \(B_i\) such that
\(c_1(B_i)\to\theta\). This is BJ Definition 4.1, not an invocation of
its subsequent varying-class closure theorem. Its defining limits are
usc, since extended-continuous FS functions are usc. Write
\(\operatorname{PSH}(L)\) for this cone with \(\theta=c_1(L)\). Within
an actual ample bundle this is the genuine ample cone in G0. We do not
assume that arbitrary nested varying-class limits flatten.

Choose and retain actual data
\[L\simeq A+\mathcal O_X(D),\quad A\text{ ample rational},\quad D\geq0\text{ rational Cartier}.\tag{1.2}\]
Here and below additive bundle notation includes the specified rational
isomorphism. Section 2 constructs such data for every big rational \(L\)
on this same \(X\), without Bertini or modifications. Define
\[\sigma=\sigma_D=\log\|s_D\|_{\rm triv},\quad
Z=\{\sigma=-\infty\}=|D|^{\rm an},\quad U=X^{\rm an}\setminus Z,\quad V_D=(A^n)>0.\tag{1.3}\]
If \(D=0\), set \(\sigma=0\) and \(Z=\varnothing\). The rational
logarithm is defined by any cleared multiple and is independent of that
multiple. Scaling acts by \(\sigma(tv)=t\sigma(v)\), including at
\(-\infty\). The reference and all operators are fixed before selecting
a target. Define the entire prescribed-reference cone
\[\mathcal P_D(L)=\{u\in\operatorname{PSH}(L):u\leq\sigma+C
\text{ on all }X^{\rm an}\text{ for some }C\in\mathbb R\}.\tag{1.4}\]

\textbf{Theorem 1 (all-big prescribed divisorial-reference equation).}
For every datum just specified:

\begin{enumerate}
\def\labelenumi{\arabic{enumi}.}
\tightlist
\item
  The map \(F_D(\psi)=\sigma+\psi\) is an order-preserving bijection
  from the actual genuine \(\operatorname{PSH}(X,A)\) onto (1.4). Its
  inverse \(Q_D\) is the unique actual \(A\)-psh extension of
  \(u-\sigma\) from \(U\); values on \(Z\) are uniquely determined, not
  arbitrarily assigned. The supremum of either potential is its value at
  \(v_0\). The normalized spaces are homeomorphic in the divisorial
  topology, and are compact. Relative envelopes are genuine and
  continuous in their quotient for continuous quotient obstacles.
\item
  On the original \(X\), the genuine local metric products on \(U\),
  extended by reference-floor nonpluripolar limits, give symmetric
  positive raw mixed measures. For \(u_i\in\mathcal E_D^1\) on the same
  datum their mass is \(V_D\); for different fixed data \(L_i=A_i+D_i\)
  their mixed mass is \((A_1\cdots A_n)\). These products equal the
  genuine ample mixed products of the unique quotients, because of an
  actual local metric isometry and removal of the pole sets. They have
  unbounded decreasing-net continuity and the stated integrable
  quotient-weighted continuity.
\item
  The actual relative energy is
  \[E_D(u)=\frac1{n+1}\sum_{j=0}^n\int Q_Du\,T_u^j\wedge T_\sigma^{n-j}
   =E_{\rm raw,A}(Q_Du),\tag{1.5}\] first for bounded quotients, then by
  reference floors. Its space
  \(\mathcal E_D^1=\{u\in\mathcal P_D(L):E_D(u)>-\infty\}\) consists of
  actual singular metrics. One has \(E_D(\sigma)=0\) and
  \(E_D(u+c)=E_D(u)+V_Dc\).
\item
  The intrinsic target condition is \[\mu(X^{\rm an})=V_D,\qquad
  D_D(\mu):=\sup_{u\in\mathcal E_D^1}
  \left(E_D(u)-\int Q_Du\,d\mu\right)<\infty.\tag{1.6}\] It is
  equivalent to integrability of every actual quotient \(Q_Du\) with
  \(u\in\mathcal E_D^1\). It implies \(\mu(Z)=0\); no support
  restriction is imposed separately. For every such target an actual
  \(u\in\mathcal E_D^1\) satisfies \(\operatorname{MA}_D(u)=\mu\), and
  all solutions differ by a real constant. The dual supremum is
  attained. Comparison and domination hold for this original cone and
  operator.
\item
  For every \(u\in\mathcal E_D^1\) and every \(f\in C^0(X^{\rm an})\)
  the reference outer energy satisfies
  \[\left.\frac d{dt}\right|_{t=0}E_D^{\downarrow}(u+tf)
   =\int f\,\operatorname{MA}_D(u),\tag{1.7}\] with both signs of \(t\).
  Its definition retains the reference even where \(u=-\infty\).
\end{enumerate}

\textbf{Theorem 2 (actual trivial-bundle zero-volume range).} For
\(L=\mathcal O_X\), the literal genuine numerical zero cone is
\(\operatorname{PSH}(0)=\mathbb R\). With reference \(0\), the raw
energy is zero and the raw operator is the zero measure. Its intrinsic
positive target class consists solely of \(\mu=0\). Every constant
solves that equation, and uniqueness modulo constants holds. There is no
division by volume and no probability normalization. This assertion is
not made for every zero-volume pseudoeffective bundle.

Theorem 1 includes every big nonample rational bundle, every fixed datum
(1.2), and every target satisfying (1.6). It does not identify different
choices of \(D\) or their masses. It concerns a prescribed divisorial
reference, rather than claiming that all of \(\operatorname{PSH}(L)\)
has been given a reference-free mass or energy. Section 11 distinguishes
this complete theorem from the full original axis.

\section{2. Actual no-Bertini data for every big rational
bundle}\label{actual-no-bertini-data-for-every-big-rational-bundle}

Use the numerical conventions in the supplied BJ source:
\(\operatorname{Psef}(X)\) is the closed cone generated by effective
real Cartier divisor classes, and \(\operatorname{Big}(X)\) is its
interior. The ample cone is open in finite-dimensional \(N^1(X)\).
Ampleness of an actual rational bundle is numerical on our projective
\(X\): a numerically ample actual representative is ample. This
algebraic fact is an external premise with its precise primary scope
recorded in Appendix A; it is not a nonample analytic grant.

Let \(\theta=c_1(L)\) be any rational big class. Fix an actual ample
integral \(H\). For sufficiently small rational \(\varepsilon>0\),
\(\theta-\varepsilon c_1(H)\) lies in \(\operatorname{Psef}(X)\). Since
\(\varepsilon c_1(H)\) is ample and that cone is open, an effective
rational Cartier divisor \(D\) can be chosen with \(c_1(D)\) close
enough to \(\theta-\varepsilon c_1(H)\) that
\[\theta-c_1(D)\in\operatorname{Amp}(X).\tag{2.1}\] Here is why the
approximation is literal. The defining closed effective cone is the
closure of finite nonnegative real combinations of effective Cartier
divisors. Each such finite combination can be approximated by
nonnegative rational coefficients. On smooth \(X\) the prime components
are Cartier, so the approximation is an actual effective rational
Cartier divisor. No infinite divisor, finitely generated section ring,
or change of field is introduced. Set \(A=L-\mathcal O_X(D)\) as an
actual rational bundle. By numerical ampleness, (2.1) makes this actual
\(A\) ample. The identity (1.2) is then literal in
\(\operatorname{Pic}(X)_{\mathbb Q}\), not merely numerical.

This proves existence for every big actual rational bundle, even with
arbitrary numerically trivial twists. It uses neither Bertini nor a
birational model. It also avoids confusing top self-intersection with
volume. Conversely, actual \(A+D\) data are big: the ample cone lies in
the interior of the effective cone (sufficiently high ample powers have
sections, and small rational perturbations remain ample), and adding an
effective class preserves that interior. The data therefore have exactly
the required big domain.

We will repeatedly use the elementary section fact that for a generating
finite family \(t_j\) of a cleared ample power,
\[\max_j\log\|t_j\|_{\rm triv}=0\quad\text{at every analytic point}.\tag{2.2}\]
At the valuation-ring center, every section is a regular local
coefficient of a frame of norm one, hence has norm at most one. A
generating section is a unit at that center and has norm one. This
includes points with nongeneric support. Global generation of
sufficiently large ample powers has precisely the hypotheses of Stacks
01Q3; those hypotheses hold on every projective \(X\) here. Products,
powers and union of the finitely many displayed section families respect
their canonical norms and actual bundles.

\section{3. Direct foundations from a single defining FS
net}\label{direct-foundations-from-a-single-defining-fs-net}

\subsection{3.1. Rational rounding without nested
closure}\label{rational-rounding-without-nested-closure}

Let \(\phi_i\downarrow\phi\) be its one defining real FS net in \(B_i\).
For a positive rational \(\delta\), round every finite coefficient
\(\lambda_{ij}\) upward to a rational number with error at most
\(m_i\delta\). The corresponding actual rational FS function
\(f_{i,\delta}\) satisfies
\[\phi_i\leq f_{i,\delta}\leq\phi_i+\delta.\tag{3.1}\] These
inequalities are non-strict and are valid also at common section zeros.
Put \(g_{i,\delta}=f_{i,\delta}+\delta\). Order pairs by increasing
\(i\), decreasing \(\delta\), and decreasing \(g\) pointwise. This set
is directed: given two pairs, take \(i_3\) above both and
\(2\delta_3\leq\min(\delta_1,\delta_2)\). Then
\[g_{i_3,\delta_3}\leq\phi_{i_3}+2\delta_3
\leq\min(\phi_{i_1}+\delta_1,\phi_{i_2}+\delta_2)
\leq\min(g_{i_1,\delta_1},g_{i_2,\delta_2}).\tag{3.2}\] The new net
decreases to \(\phi\) everywhere: at a finite value choose an original
index and then a sufficiently small \(\delta\); at \(-\infty\) use any
finite threshold. The index map is increasing and final onto the
original index set, and \(\delta\to0\), so its actual bundle classes
still tend to \(\theta\). If \(\phi(v_0)\geq0\), the new approximants
retain nonnegative generic values. This repairs rational approximation
directly, without BJ Lemma 4.6.

\subsection{3.2. Embedding every defining numerical function into an
actual ample
cone}\label{embedding-every-defining-numerical-function-into-an-actual-ample-cone}

Choose an actual rational ample \(H\) with \(c_1(H)-\theta\) ample.
Eventually \(H-B_i\) is an actual ample rational bundle. After clearing
degrees, multiply every displayed \(B_i\)-section by a generating family
of a sufficiently divisible power of \(H-B_i\). Formula (2.2) makes the
new actual \(H\)-section maximum exactly \(\phi_i\), including all its
poles. Thus the decreasing tail is a defining FS net on the single
actual ample \(H\), and \[\phi\in\operatorname{PSH}(X,H).\tag{3.3}\]
Finite lists of numerical functions can be placed in one common \(H\).
The ample grant now gives finiteness on every divisorial point,
\(\sup\phi=\phi(v_0)\), all-point divisorial determination, and the
implication from a divisorial inequality to an arbitrary usc obstacle.
These consequences are valid for our defining numerical functions by
(3.3); no nonample closure theorem is being assumed.

\subsection{3.3. Homogeneous projection by a finite
polytope}\label{homogeneous-projection-by-a-finite-polytope}

For \(\phi(v_0)\geq0\), set
\[\widehat\phi(v)=\inf_{t>0}\frac{\phi(tv)}t.\tag{3.4}\] This operation
preserves order, is homogeneous, and commutes with decreasing infima:
\(\inf_i\inf_t=\inf_t\inf_i\), also in the extended real sense. For a
rational FS function (1.1), put \(x_j=\log\|s_j\|(v)\), and write
\(r=1/t\). Then \[m\widehat\phi(v)=\inf_{r\geq0}\max_j(x_j+r\lambda_j)
=\max_{w\in W}\sum_jw_jx_j,\tag{3.5}\] where \(W\) is the finite vertex
set of
\[P=\{w_j\geq0,\ \sum_jw_j=1,\ \sum_jw_j\lambda_j\geq0\}.\tag{3.6}\] The
infimum for \(r>0\) equals that for \(r\geq0\) by continuity at zero.
For finite \(x_j\), every feasible weighted expression is at most the
maximum for every \(r\). Conversely a minimum of the convex piecewise
affine function in \(r\) occurs at zero, at a finite breakpoint, or
asymptotically at infinity. At zero an active nonnegative slope gives an
eligible unit vertex. At a breakpoint a zero active slope or two
opposite active slopes with their zero weighted average gives a vertex
attaining the same value. The asymptotic finite case has largest slope
zero and is attained by an active zero-slope unit vertex. These cases
prove equality. The vertices are the eligible unit vectors and the zero
averages of two opposite slopes; their coefficients are rational.

Choose a common denominator \(b\) of the weights. The monomials
\[S_w=\prod_j s_j^{bw_j}\in H^0(X,mbB),\qquad w\in W,\tag{3.7}\] are
actual nonzero sections on integral \(X\). Their homogeneous FS maximum,
divided by \(mb\), is (3.5). Zero exponents mean omission of the
section, so no \(0\cdot(-\infty)\) ambiguity occurs. If some
\(x_j=-\infty\), replace them by \(\max(x_j,-N)\) in (3.5), commute the
infima over \(N\) and \(r\), and then pass the fixed finite maximum of
weighted sums to its decreasing limit. This proves the same formula at
every base-locus point, including the case where it equals \(-\infty\).

Apply this calculation to the rational net of Section 3.1. The
homogeneous actual FS functions decrease to \(\widehat\phi\), have
generic value zero, and have actual classes tending to \(\theta\). They
therefore prove \(\widehat\phi\in\operatorname{PSH}_{\rm def}(\theta)\)
by the original single-net definition. No varying-class closure or
printed homogeneous approximation theorem is needed. If \(h\) is
homogeneous, \(\widehat h=h\).

\subsection{3.4. All-point homogeneous
decomposition}\label{all-point-homogeneous-decomposition}

Let \(C=\phi(v_0)\) and, for \(\lambda<C\), put
\(h_\lambda=\widehat{\phi-\lambda}\). It is a genuine homogeneous
numerical \(\theta\)-psh function with generic value zero, and
\(h_\lambda+\lambda\leq\phi\). At every actual point \(v\),
\[\phi(v)=\sup_{\lambda<C}\{h_\lambda(v)+\lambda\}.\tag{3.8}\] To prove
the reverse inequality, consider \(g(t)=\phi(tv)\) on \((0,\infty)\).
Each defining FS function on this ray is convex, decreasing, and bounded
above by its generic value. The decreasing limit has the same
properties. If it is \(-\infty\) at one point, convexity using a larger
point with a finite uniform upper bound forces \(-\infty\) at every
smaller point, and monotonicity forces it at every larger point. Thus
the ray is either identically \(-\infty\) or finite convex. The first
case proves (3.8) immediately.

In the finite case choose a supporting line \(at+\lambda_0\) at \(t=1\).
Its slope is nonpositive and its intercept is at most \(C\), by taking
\(t\downarrow0\) in the inequality \(at+\lambda_0\leq g(t)\leq C\). If
\(\lambda_0<C\), (3.4) gives \(h_{\lambda_0}(v)\geq a\), so equality at
\(t=1\) proves (3.8). If \(\lambda_0=C\), lower its intercept to
\(C-\epsilon\): the same line inequality gives
\(h_{C-\epsilon}(v)\geq a\), and the right side of (3.8) is at least
\(g(1)-\epsilon\). Let \(\epsilon\downarrow0\). This proves (3.8) on all
support types and pole points without BJ 6.24.

\section{4. The entire dominated cone and its unique actual
quotient}\label{the-entire-dominated-cone-and-its-unique-actual-quotient}

\subsection{4.1. Forward membership, including common
poles}\label{forward-membership-including-common-poles}

G0 supplies a decreasing net of finite continuous actual \(A\)-FS
functions \(\psi_j\) for every \(\psi\in\operatorname{PSH}(X,A)\). Clear
the degree of \(\psi_j\) and the denominators of \(A,D,L\) to one
integer \(p\). Multiply the displayed \(pA\)-sections by the canonical
section of \(pD\). Their actual \(pL\)-section maximum is exactly
\[\sigma+\psi_j.\tag{4.1}\] The net decreases on every analytic point to
\(\sigma+\psi\). On \(Z\) all these maxima are exactly \(-\infty\). It
proves genuine membership by the single defining \(L\)-net, and
\(\sigma+\psi\leq\sigma+\sup\psi\). Thus \(F_D\) maps into the full cone
(1.4).

For completeness, finite fixed-bundle flattening, when needed solely in
this ample cone, is legitimate. Suppose \(\psi_i\downarrow\psi\), and
each \(\psi_i\) is a decreasing limit of finite continuous \(A\)-FS
functions \(f_{ij}\). Order triples \((i,j,\epsilon)\) by increasing
\(i\) and decreasing \(f_{ij}+\epsilon\). For two triples take a larger
outer index and a strictly smaller positive \(\epsilon\). The limiting
lower function plus that \(\epsilon\) is strictly below the minimum of
the two finite upper functions everywhere, even where the limit is
\(-\infty\). The continuous exponential differences decrease to a
strictly negative function. Their negative open sets cover compact
\(X^{\rm an}\); a finite subcover and inner directedness give one lower
finite approximant globally. This makes the triples directed and their
FS functions decrease to \(\psi\) (use arbitrary finite thresholds at
\(-\infty\)). This restricted argument has finite upper inner functions.
We never apply its strict step to (4.1) at their common poles. G0
already supplies the required finite fixed-\(A\) approximation; this
argument explains why that use does not import the failed general lemma.

\subsection{4.2. Homogeneous division using compactness, not divided-net
monotonicity}\label{homogeneous-division-using-compactness-not-divided-net-monotonicity}

Let \(h\) be homogeneous, genuine \(L\)-psh, and \(h\leq\sigma\). If
\(D=0\), \(L=A\) and the assertion is the ample identification. Assume
\(D\ne0\). Section 3.3 gives decreasing homogeneous actual rational FS
\(h_i\) in \(B_i\) with \(c_1(B_i)\to c_1(L)\) and \(h_i(v_0)=0\). The
set \[K=\{\sigma=-1\}\tag{4.2}\] is compact: \(\sigma\) is continuous
into \(\mathbb R\cup\{-\infty\}\) and \(\{-1\}\) is closed. It is
nonempty, since a component valuation of \(D\) can be scaled to it. For
every rational \(t\in(0,1)\), the open sets \(\{h_i<-t\}\cap K\)
increase and cover \(K\), since \(h\leq-1\) there. Compactness yields an
entire tail with \(h_i<-t\) on \(K\).

Use the directed cofinal pairs \((i,t)\) satisfying this condition,
ordered by increasing \(i\) and \(t\), with \(t\to1\). Homogeneity gives
\(h_i(v)\leq t\sigma(v)\) on each divisorial \(v\): scale \(v\) to \(K\)
if \(\sigma(v)<0\), and use \(h_i\leq0\) if \(\sigma(v)=0\). A fixed
actual ample embedding for each \(h_i\) and its usc obstacle gives
\[h_i\leq t\sigma\quad\text{on all }X^{\rm an}.\tag{4.3}\] This
application needs no ample membership of \(t\sigma\): it is an usc
obstacle.

Apply literal section division to (4.3). After clearing degrees, each
displayed section of \(pB_i\) vanishes along every prime component of
\(ptD\) to at least its specified multiplicity, as seen by its prime
divisor valuation. Normality of smooth \(X\) implies that division by
\(s_{ptD}\) gives an actual global section of \(p(B_i-tD)\).
Equivalently, in a normal trivializing affine ring a rational
coefficient with no pole at any height-one prime is regular; this is
exactly the normality input in BJ Lemma 2.23 and its preserved proof.
Therefore \[r_{i,t}=h_i-t\sigma\quad\text{on }U\tag{4.4}\] has a literal
homogeneous FS extension on \(B_i-tD\), given by these quotient
sections, and generic value zero. Formula (4.4) is not an extended
subtraction on \(Z\).

The divided functions need not decrease. Instead adjoin rational
\(\eta>0\) and retain triples with
\[c=1+\eta,\qquad cA-(B_i-tD)\text{ actual ample}.\tag{4.5}\] Order them
by increasing \(i,t\) and decreasing \(\eta\). They are directed and
cofinal with \(t\to1,\eta\to0\): for any smaller positive \(\eta\), take
\(t\) close enough to one and \(i\) sufficiently late that \(B_i-tD\) is
numerically close enough to \(A\) for (4.5), using openness of the ample
cone. Increasing \(i\) again retains the compact \(K\) condition. After
rational scaling and multiplication by a generating family of a cleared
power of the ample difference (4.5), (2.2) realizes
\[a_{i,t,\eta}=c^{-1}r_{i,t}\in\operatorname{PSH}(X,A),\qquad a_{i,t,\eta}(v_0)=0.\tag{4.6}\]
These are actual normalized \(A\)-FS functions.

G0 compactness gives a subnet converging to an actual
\(r\in\operatorname{PSH}(X,A)\) in the divisorial topology. At every
divisorial point, the values in (4.6) tend explicitly to
\(h(v)-\sigma(v)\): \(h_i(v)\to h(v)\), \(t\to1\), \(c\to1\), and both
\(h(v)\) and \(\sigma(v)\) are finite there. Hence \(r=h-\sigma\) on
divisorial points. By Section 4.1, \(\sigma+r\) is genuine \(L\)-psh.
Put it and \(h\) into one actual ample \(H\) by Section 3.2. Ample
divisorial determination gives
\[h=\sigma+r\quad\text{at every analytic point}.\tag{4.7}\] In
particular \(h=-\infty\) on \(Z\), as required. Divisorial determination
in the actual ample \(A\) proves uniqueness of \(r\).

\subsection{4.3. General dominated
functions}\label{general-dominated-functions}

Let \(u\in\mathcal P_D(L)\), so \(u\leq\sigma+C_0\), and put
\(C=u(v_0)\). Section 3.4 gives
\(u=\sup_{\lambda<C}(h_\lambda+\lambda)\) at all points. At a divisorial
point, evaluate \(h_\lambda+\lambda\leq u\leq\sigma+C_0\) on \(tv\),
divide by \(t\), and let \(t\to\infty\). This gives
\(h_\lambda\leq\sigma\) divisorially, then everywhere by the ample
embedding and usc obstacle principle. Section 4.2 gives actual
\(r_\lambda\in\operatorname{PSH}(X,A)\) with
\[h_\lambda=\sigma+r_\lambda,\qquad r_\lambda(v_0)=0.\tag{4.8}\] Thus
\(r_\lambda+\lambda\leq C\) everywhere.

We reconstruct its ample supremum without assuming a nonample supremum
theorem. For finite nonempty \(J\subset\{\lambda<C\}\), let
\[b_J=\max_{\lambda\in J}(r_\lambda+\lambda),\qquad c_J=\max J.\tag{4.9}\]
They are genuine actual \(A\)-psh, \(b_J(v_0)=c_J\), and
\(c_J\uparrow C\). The normalized \(b_J-c_J\) have a compact convergent
subnet by G0. On each divisorial point, \(b_J(v)\) increases to a finite
value: it is bounded above by \(C\) and below by any one fixed
\(r_\lambda(v)+\lambda\). The subnet therefore gives an actual
\(r\in\operatorname{PSH}(X,A)\) with
\[r(v)=\sup_{\lambda<C}(r_\lambda(v)+\lambda)\quad(v\in X^{\rm div}),\qquad r(v_0)=C.\tag{4.10}\]
For each \(\lambda\), ample divisorial domination yields
\(r_\lambda+\lambda\leq r\) at every point. Taking the usc
regularization of the supremum gives
\(\sup^*(r_\lambda+\lambda)\leq r\). The reverse inequality follows from
(4.10) and the ample usc obstacle principle. Hence \(r\) is exactly that
usc supremum. On divisorial points (3.8) gives \(u=\sigma+r\). Both
sides are genuine \(L\)-psh and lie in a common ample \(H\). Ample
divisorial determination proves equality on all \(X^{\rm an}\). This
proves surjectivity of \(F_D\) for the entire cone (1.4), not merely for
a selected FS subcone.

If \(F_D(\psi_1)=F_D(\psi_2)\), subtract the finite \(\sigma\) on
divisorial points and use ample determination to obtain
\(\psi_1=\psi_2\) everywhere. Likewise \(F_D(\psi_1)\leq F_D(\psi_2)\)
iff \(\psi_1\leq\psi_2\). In the forward implication, the pole values
cause no subtraction: divisorial inequalities suffice. This proves the
claimed order bijection and the unique all-point quotient \(Q_D\).

\section{5. Normalization, compactness and genuine relative
envelopes}\label{normalization-compactness-and-genuine-relative-envelopes}

The generic normalization is exact:
\[u(v_0)=Q_Du(v_0),\qquad\sup u=u(v_0),\qquad\sup Q_Du=Q_Du(v_0).\tag{5.1}\]
The equality of suprema for \(u\) uses Section 3.2 and G0, rather than a
heuristic density argument. On every divisorial coordinate \(F_D\) is
addition of the fixed finite number \(\sigma(v)\), and \(Q_D\) is
subtraction of that number. Thus the maps are inverse homeomorphisms in
the stated divisorial topology. G0 compactness of the normalized actual
ample cone proves compactness of the original normalized cone (1.4).
These are actual functions; there is no completion object or
modification.

A family in (1.4) uniformly bounded above has uniformly bounded generic
values, hence uniformly bounded quotients, by (5.1). Conversely
uniformly bounded quotients imply a common bound by \(\sigma+C\). Finite
maxima and reference-relative convex combinations commute with \(F_D\);
at \(Z\) all sums take value \(-\infty\), and elsewhere this is ordinary
arithmetic. If \(u_i\downarrow u\) in this cone and the endpoint is
generically finite, their quotients decrease: order reflection gives the
order, G0 ample decreasing closure gives an actual limit, and ample
determination identifies its image with \(u\) on every point. This is a
proved fixed-reference closure, not arbitrary varying-class closure.

For any usc \(b\) bounded above, the exact envelope correspondence is
\[P_D(\sigma+b)=F_D(P_A(b)).\tag{5.2}\] Indeed \(F_D(\psi)\leq\sigma+b\)
implies \(\psi\leq b\) on divisorial points and then everywhere by the
ample usc obstacle principle. Conversely \(\psi\leq b\) implies the
original inequality, including at \(Z\). The envelopes and their usc
regularizations consequently match; for a uniformly bounded family
\(\psi_\alpha\),
\[\sup^*_\alpha F_D(\psi_\alpha)=F_D(\sup^*_\alpha\psi_\alpha).\tag{5.3}\]
On \(U\), \(\sigma\) is finite continuous, so usc regularization
commutes with its addition. At \(Z\), the common upper bound on
quotients and \(\sigma\to-\infty\) force both sides to be \(-\infty\).
Thus G0 EP supplies genuine relative EP at every point. For continuous
\(b\) the quotient \(P_A(b)\) is continuous, and (5.2) is a continuous
singular metric in the specified reference sense. It is not a globally
finite metric on a nonample \(L\).

\section{6. Genuine metric intersections and reference-floor
extension}\label{genuine-metric-intersections-and-reference-floor-extension}

\subsection{6.1. The actual local metric
isometry}\label{the-actual-local-metric-isometry}

Potentials use \(\|e\|_\psi=e^{-\psi}\|e\|_{\rm triv}\). Clear rational
powers and write in an actual local frame \(s_D=f e_D\), \(e_L=e_Ae_D\).
For \(u=\sigma+\psi\), the local logarithmic metric weight is
\[h_L=-\log\|e_L\|_u
 =\psi-\log\|e_A\|_{\rm triv}+\log|f|
 =h_A+\log|f|.\tag{6.1}\] On \(U\), the regular nowhere vanishing
section \(s_D\) identifies \(A|_U\) with \(L|_U\). The frame
\(e_As_D=f e_L\) has weight \(h_L-\log|f|=h_A\). Thus the two actual
metrized bundles are locally isometric. The sign in (6.1) is plus. A
claim with the opposite sign would describe the wrong metric.

The public CLD unit states that the log norm of an invertible analytic
function is smooth with zero \(d'd''\); its continuous Bedford--Taylor
product is invariant under these harmonic frame terms. Consequently
actual continuous quotient metrics have on \(U\) exactly the genuine
original \(L\) curvature products of their \(A\) metrics, not an inverse
measure assigned after choosing a target. For actual \(A\)-FS metrics
the regularized maxima of BE 7.12 give smooth semipositive approximants;
BE 7.14 applies to semiample bundles over complete fields including the
trivial valuation. BE 7.8, which assumes a nontrivial valuation, is not
invoked. BE 8.3--8.5 give the genuine positive fundamental-cycle
products and uniform continuity. G0 identifies these continuous products
with its genuine intersection operator. All these hypotheses hold for
the actual ample \(A\).

For bounded, possibly discontinuous actual
\(\psi_i\in\operatorname{PSH}(A_i)\), take decreasing finite continuous
FS nets supplied by G0. On \(U=X^{\rm an}\setminus\bigcup Z_i\),
transport those metrics by (6.1) first. The local products have the
limit \[T_{u_1}\wedge\cdots\wedge T_{u_n}|_U
 =\bigwedge_{i=1}^n(c_1(A_i)+dd^c\psi_i)|_U.\tag{6.2}\] For a continuous
test with compact support in \(U\), extend it by zero to \(X^{\rm an}\)
and apply G0 mixed decreasing-net convergence. This proves independence
of the chosen global approximants and actual local original-carrier
agreement. It constructs the bounded metric intersection in this cone;
it does not purport to prove a local Bedford--Taylor theorem for
unrelated arbitrary local psh inputs.

\subsection{6.2. A polar test, with the capacity exponent made
explicit}\label{a-polar-test-with-the-capacity-exponent-made-explicit}

We need to justify removal of \(Z_i\) for unbounded finite-energy
products and targets. The following proof uses only actual ample cones.
For any actual ample \(A\), choose \(q\) and a nonzero
\(s\in H^0(X,qA)\) vanishing along \(|D|\). Such a section exists by
global generation of \(\mathcal I_{|D|}\otimes qA\) for large \(q\), and
is nonzero at some point outside \(|D|\). Set
\(a=q^{-1}\log\|s\|_{\rm triv}\), with \(a(v_0)=0\). Its floors
\(\max(a,-j)\) are finite actual \(A\)-FS functions, obtained by adding
a sufficiently divisible generating section family of constant weight
\(-j\). Their decreasing limit is genuine \(A\)-psh by G0. This
construction imposes no extra condition on \(D\).

Normalize ample Monge--Ampere measures to probability only inside this
positive-mass estimate and write
\[\operatorname{Cap}_A(B)=\sup_{-1\leq b\leq0}\operatorname{MA}_{\rm BJ,A}(b)(B).\]
For every normalized ample \(a\), bounded \(-1\leq b\leq0\) satisfies
\(\int(-a)\operatorname{MA}_{\rm BJ}(b)\leq n\). At finite stages this
follows by telescoping the difference from
\(\operatorname{MA}(0)=\delta_{v_0}\), integration by parts, and
\(-1\leq b\leq0\): each signed difference of two positive mixed terms of
normalized mass one is bounded in absolute value by one. Floor
approximation and monotone convergence give the assertion for unbounded
\(a\). This is the ample content of the preserved BJ 7.32 argument.
Therefore \[\operatorname{Cap}_A\{a\leq-t\}\leq n/t.\tag{6.3}\]

For a normalized ample potential \(b\), let \(b_s=\max(b,-s)\) and
\(\mu_s=\operatorname{MA}_{\rm BJ}(b_s)\). Energy comparison gives
\(J(b_s)\leq\int(-b_s)d\mu_s\). Locality and equal mass imply
\(\mu_s\{b\leq-t\}=\mu_t\{b\leq-t\}\) for \(1\leq t\leq s\). Since
\(t^{-1}b_t\) is actual \(A\)-psh and lies in \([-1,0]\), expansion of
its positive mixed product gives
\(\operatorname{MA}(t^{-1}b_t)\geq t^{-n}\mu_t\). Hence
\[J(b)\leq1+\int_1^\infty t^n\operatorname{Cap}_A\{b\leq-t\}\,dt.\tag{6.4}\]
The weight \(t^n\) must be retained; some supplied renderer lines omit
it. We use the strictly sufficient exponent \(0<\alpha<1/(n+1)\),
avoiding reliance on the more liberal printed exponent in a corollary.

The convex increasing function \(\chi(x)=1-(1-x)^\alpha\), \(x\leq0\),
has \(0\leq\chi'\leq1\). Its composition \(b=\chi(a)\) is genuine
\(A\)-psh: approximate \(\chi\) decreasingly by convex piecewise affine
chord majorants on expanding compact intervals, with constant left tails
and slopes in \([0,1]\), use ample convexity, and then G0 fixed-ample
decreasing closure. It is normalized and \(b=-\infty\) on \(Z\). By
(6.3), \[\operatorname{Cap}_A\{b\leq-t\}=O(t^{-1/\alpha}).\] Since
\(n-1/\alpha<-1\), (6.4) gives \(b\in\mathcal E^1(X,A)\). Thus \(Z\) is
contained in the pole set of an actual ample finite-energy potential.
The same proof applies to any pole set defined by a numerical potential:
Section 3.2 first embeds it into some actual ample \(H\); scaling into
\(A\) can be done directly on its defining sections after adding
generating sections of an ample difference. All steps occur inside fixed
ample cones after that embedding.

G0 mixed integrability now proves that every finite-energy ample mixed
measure gives zero mass to these polar sets: it integrates the just
constructed finite-energy test. Likewise any target that integrates
every actual ample \(\mathcal E^1\) potential gives zero mass to \(Z\).
No unproved negligible-set theorem, Choquet extraction, or countability
assumption is needed.

\subsection{6.3. Reference floors and raw
mass}\label{reference-floors-and-raw-mass}

For arbitrary \(u=F_D(\psi)\) define genuine reference floors
\[\psi_j=\max(\psi,-j),\quad
u_j=F_D(\psi_j)=\max(u,\sigma-j),\quad
B_j=U\cap\{\psi>-j\}.\tag{6.5}\] The bounded local metric product of
\(\nu_j\) exists by Section 6.1. Its restriction to \(B_j\) increases
with \(j\): on overlapping strict winning sets, G0 genuine locality in
the quotient gives equality, and \(B_j\) increases. The sets are Borel;
they need not be ordinary open sets. Extend the restricted measures by
zero and set
\[\operatorname{MA}_D(u)=\lim_{j\to\infty}\mathbf1_{B_j}T_{\nu_j}^n.\tag{6.6}\]
This construction precedes any target measure. Each positive measure has
mass at most \(V_D\). Their increasing limit is a finite Radon measure:
for any Borel set, approximate one late finite measure from inside by a
compact set; monotonicity and finite limiting mass prove inner
regularity of the limit. On a compact Hausdorff space finite inner
regularity also gives the requisite outer regularity by complements.
Thus (6.6) is the genuine reference-floor nonpluripolar extension, of
mass at most \(V_D\) for general \(\psi\).

If \(\psi\in\mathcal E^1(X,A)\), quotient locality gives
\[\mathbf1_{B_j}T_{\nu_j}^n
 =\mathbf1_{B_j}\operatorname{MA}_{\rm raw,A}(\psi).\] The latter
measure integrates \(\psi\) and gives zero mass to \(Z\) by Section 6.2.
The union of \(B_j\) has full measure. Consequently
\[\operatorname{MA}_D(F_D\psi)=\operatorname{MA}_{\rm raw,A}(\psi),\quad
\operatorname{MA}_D(F_D\psi)(X^{\rm an})=V_D.\tag{6.7}\] This equality
has been derived from actual local metric products, locality, and a
polar test. It is not an alternative inverse measure definition.

For different data and finite-energy \(\psi_i\), use the common
complement of their \(Z_i\), the simultaneous quotient floors, and
strict winning sets \(\bigcap_i\{\psi_i>-j\}\). Mixed locality follows
from finite algebraic polarization of the diagonal locality identities.
The limiting product is the positive genuine mixed product of the
\(\psi_i\), of mass \((A_1\cdots A_n)\), and gives no mass to their pole
sets. Polarization has alternating coefficients; positivity here follows
from the actual mixed product, not from calling the polarization
coefficients positive.

For decreasing genuine relative nets \(u_{i\beta}\downarrow u_i\) with
finite-energy endpoints, Section 5 proves
\(Q_{D_i}u_{i\beta}\downarrow Q_{D_i}u_i\) everywhere. G0 and (6.7)
prove weak convergence of all unbounded mixed measures with their fixed
raw intersection masses. If \(\psi_{0\beta}\downarrow\psi_0\) is a
genuine ample test net and \(\psi_0\) is integrable for the limiting
mixed measure, the signed weighted pairing is
\[\int\psi_0\,\bigwedge_{i=1}^nT_{u_i}
=(c_1(A_0),\psi_0)\prod_{i=1}^n(c_1(A_i),\psi_i)
 -(c_1(A_0),0)\prod_{i=1}^n(c_1(A_i),\psi_i).\tag{6.8}\] G0 decreasing
pairing continuity, the fixed positive masses, and the integrable-weight
convergence give weak convergence of the weighted measures. Formula
(6.8) keeps the correct subtraction order. Finite-energy quotient tests
are integrable by G0 mixed integrability. Differences are handled by
their two terms, not by declaring a difference of decreasing nets
decreasing. We make no extra assertion that \(\sigma\) itself is
integrable against every admissible target or product; the meaningful
relative weight is its unique quotient.

\section{7. Actual energy and the intrinsic target
class}\label{actual-energy-and-the-intrinsic-target-class}

For bounded quotients, define energy by the genuine products already
constructed in Section 6: \[E_D(u)=\frac1{n+1}\sum_{j=0}^n\int Q_Du\,
T_u^j\wedge T_\sigma^{n-j}.\tag{7.1}\] Every integral is a finite
integral on the original \(X\). Equation (6.7) identifies each mixed
factor with the ample factor of \(\psi=Q_Du\) and reference \(0\). The
usual ample intersection energy formula in G0 therefore gives
\(E_D(u)=E_{\rm raw,A}(\psi)\). This equality has its local geometric
justification in Section 6.1 and its pole removal justification in
Section 6.2; it is not an abstract assignment to a lift. For arbitrary
\(u\) define \[E_D(u)=\lim_{j\to\infty}E_D(\max(u,\sigma-j))
=E_{\rm raw,A}(Q_Du)\in\mathbb R\cup\{-\infty\}.\tag{7.2}\] Ample
monotonicity and decreasing continuity prove independence of the chosen
cofinal real reference floors. It follows that
\[\mathcal E_D^1=F_D(\mathcal E^1(X,A)),\quad
E_D(\sigma)=0,\quad E_D(u+c)=E_D(u)+V_Dc.\tag{7.3}\] This supplies the
actual finite-energy space, reference and mass explicitly. Concavity,
monotonicity, and weak upper semicontinuity follow from the order
homeomorphism and the actual G0 energy. All finite mixed energies and
weighted identities retain the raw masses from Section 6. For normalized
\(u\), define \(J_D(u)=-V_D^{-1}E_D(u)\), equal to the normalized ample
\(J\).

The quotient is intrinsic to \((L,\sigma)\): it is the unique actual psh
extension of the scalar difference on the reference complement.
Different choices of that extension are impossible by Section 4.
Intrinsic dual energy is exactly (1.6). Put \(\bar\mu=\mu/V_D\) only
when \(V_D>0\). The identities are
\[D_D(\mu)=V_D E_{\rm BJ,A}^{\vee}(\bar\mu),\quad
\operatorname{MA}_D(u)=V_D\operatorname{MA}_{\rm BJ,A}(Q_Du).\tag{7.4}\]
The variational objective is constant invariant precisely when
\(\mu(X^{\rm an})=V_D\). If the mass is different, translating \(u\)
makes the objective unbounded above. Thus the mass is part of the target
condition, not a normalization convention applied afterwards.

G0 on the actual ample \(A\) proves that finite (7.4) is equivalent to
integrability of every actual \(\mathcal E^1(X,A)\) potential. The
bijection (7.3) makes this exactly the all-\(\mathcal E_D^1\) quotient
integrability assertion of Theorem 1. For clarity, the nontrivial
direction in that ample criterion uses bounded-\(J\) integrability,
finite positive convex sums with unused coefficient on zero, fixed-ample
decreasing closure, and the ample scaling estimate for \(J\); it does
not select a solution or restrict measure support. Those are all
actual-ample statements explicitly contained in G0 and are applied only
to \(A\). The polar test in Section 6.2 then implies \(\mu(Z)=0\) as a
consequence of the intrinsic condition. Absolute integrability of \(u\)
or of \(\sigma\) is not added: subtracting the fixed reference is the
correct relative quantity. The condition is defined before the inverse
equation and cannot be satisfied merely by declaring a measure to be an
inverse image.

\section{8. Every continuous original variation and attained
equation}\label{every-continuous-original-variation-and-attained-equation}

\subsection{8.1. Outer energy retaining the
reference}\label{outer-energy-retaining-the-reference}

For a reference-dominated upper bounded original expression \(u+tf\),
retain its prescribed all-point quotient \(Q_Du+tf\). Define the
reference outer energy using continuous quotient majorants:
\[E_D^{\downarrow}(u+tf)
=\inf_{b\in C^0(X^{\rm an}),\ b\geq Q_Du+tf}
E_D^{\downarrow}(\sigma+b)
=E_{\rm raw,A}^{\downarrow}(Q_Du+tf).\tag{8.1}\] Here
\(E_D^{\downarrow}(\sigma+b)=E_{\rm raw,A}(P_A(b))\) by (5.2).
Continuous majorants form a decreasing directed approximation net to an
upper bounded usc quotient on compact Hausdorff \(X^{\rm an}\), and G0
outer energy gives the last equality. Equivalently it is the energy of
the exact relative envelope of that original obstacle. At \(Z\) the
original scalar obstacle is \(-\infty\), so literal subtraction would
lose values; (8.1) records the uniquely retained quotient rather than
choosing them arbitrarily. If a continuous \(b\) only satisfies
\(\sigma+b\geq u+tf\) originally, its divisorial inequality implies
\(b\geq Q_Du+tf\) everywhere by the ample usc obstacle principle. Thus
the same infimum is obtained from the original inequality itself. This
makes (8.1) an intrinsic reference construction.

G0's actual-ample two-sided outer derivative applies to
\(\psi=Q_Du\in\mathcal E^1(X,A)\) and every continuous original \(f\).
Combining (8.1) and (6.7) proves (1.7). The extension from PL tests does
not lose a sign: if \(\|f-f_m\|_\infty\leq\epsilon_m\), then
\[\psi+tf_m-|t|\epsilon_m\leq\psi+tf\leq\psi+tf_m+|t|\epsilon_m.\]
Monotonicity and the raw translation factor bound the energy error by
\(V_D|t|\epsilon_m\). Apply the ample uniform PL derivative first,
divide by positive or negative \(t\) with the proper reversed
inequality, and let \(m\to\infty\). Both derivatives equal
\(\int f\,\operatorname{MA}_D(u)\). All continuous variations are on the
original carrier; none is required to descend from a skeleton or
modification.

\subsection{8.2. Coercivity and an actual
maximizer}\label{coercivity-and-an-actual-maximizer}

Let \(\mu\) satisfy (1.6), write \(d=D_D(\mu)/V_D\), and take normalized
actual \(u_j\in\mathcal E_D^1\) whose raw objectives tend to
\(D_D(\mu)\). Their actual quotients \(\psi_j=Q_Du_j\) satisfy
\(\sup\psi_j=0\). The ample strict-concavity coercivity estimate in G0
says \[J(\psi_j)\leq C_n\left(J_{\bar\mu}(\psi_j)+d\right),\qquad
J_{\bar\mu}(\psi)=d-E_{\rm BJ,A}(\psi)+\int\psi\,d\bar\mu.\tag{8.2}\]
The second term tends to zero, so these \(J\)'s are bounded. The
normalized compactness proved in Section 5 supplies a subnet converging
to an actual normalized \(u\in\mathcal P_D(L)\), with actual quotient
\(\psi\). Lower semicontinuity of \(J\) gives \(J(\psi)<\infty\), hence
\(u\in\mathcal E_D^1\). This is net compactness on potentially
nonmetrizable spaces; the countability of a maximizing sequence does not
assert a sequentially compact analytic space or make the field
countable.

The ample bounded-\(J\) target-integral continuity in G0 yields
\(\int\psi_jd\mu\to\int\psi d\mu\) on that subnet. Upper semicontinuity
of raw ample energy, transported by (7.2), therefore gives
\[E_D(u)-\int Q_Du\,d\mu\geq D_D(\mu).\tag{8.3}\] The reverse inequality
is its defining supremum. This proves actual attainment by a genuine
original function. We have not replaced this step by convergence of
output measures, an abstract metric completion, or an arbitrary lift.

\subsection{8.3. Both Euler inequalities and genuine
existence}\label{both-euler-inequalities-and-genuine-existence}

For \(f\in C^0(X^{\rm an})\) and real \(t\), the relative envelope of
\(u+tf\) corresponds to \(P_A(\psi+tf)\), lies between
\(\psi\pm|t|\|f\|_\infty\) in its quotient, and is actual finite energy.
Feasibility and maximality (8.3) give
\[E_D^{\downarrow}(u+tf)-\int(\psi+tf)\,d\mu
\leq E_D(u)-\int\psi\,d\mu.\tag{8.4}\] Indeed the envelope's objective
is at least the left side, since its quotient is no larger than the
obstacle; it is at most the attained maximum. Differentiating (8.4) for
\(t>0\) gives \(\int f\operatorname{MA}_D(u)\leq\int f d\mu\). For
\(t<0\), division reverses the inequality and gives the opposite
inequality. Consequently the integrals agree for every continuous \(f\).
Riesz uniqueness on the compact original analytic space gives
\[\operatorname{MA}_D(u)=\mu.\tag{8.5}\] The operator used here was the
original local metric and reference-floor operator of Section 6,
constructed before selecting \(\mu\). Existence thus proves a genuine
equation. Conversely the ample concavity inequality and (7.4) show that
every finite-energy solution has finite dual energy and attains the dual
supremum, so (1.6) is exactly the solved measure class.

\section{9. Comparison, domination and uniqueness on the original
carrier}\label{comparison-domination-and-uniqueness-on-the-original-carrier}

For \(u=F_D(\psi)\) and \(v=F_D(\chi)\), on \(U\) one has \(u<v\) iff
\(\psi<\chi\). On \(Z\) both original scalars are \(-\infty\); the
operator and every target in (1.6) give that set zero mass. G0
comparison in the actual ample quotient therefore gives precisely
\[\int_{\{u<v\}}\operatorname{MA}_D(v)
\leq\int_{\{u<v\}}\operatorname{MA}_D(u).\tag{9.1}\] Strict winning
locality follows the same way and was already used to construct
reference-floor products. Mixed locality follows by the finite
polarization justified in Section 6. Domination follows from the ample
domination principle and order reflection: if \(u\leq v\) almost
everywhere for \(\operatorname{MA}_D(v)\), the quotient inequality holds
almost everywhere off \(Z\), hence everywhere for the actual ample
quotients, and \(u\leq v\) originally at all points.

Define the original finite-energy distance functional by its actual
relative differences: \[I_D(u,v)=\int(\psi-\chi)
\left(\operatorname{MA}_D(v)-\operatorname{MA}_D(u)\right)
=V_DI_A(\psi,\chi)\geq0.\tag{9.2}\] The quotient differences are
integrable by G0; this does not subtract two original \(-\infty\)
values. If the two raw output measures coincide, \(I_D=0\). Ample
uniqueness gives \(\psi-\chi=c\) at every point, and hence \(u=v+c\) on
all of \(X^{\rm an}\). Conversely translating a metric preserves its
products. Connected integral \(X\) yields one constant. This proves the
uniqueness claim in the actual cone, including the retained reference
values.

\section{10. Numerical zero foundations and the separate zero-mass
range}\label{numerical-zero-foundations-and-the-separate-zero-mass-range}

Let \(\phi\in\operatorname{PSH}_{\rm def}(0)\) have its one defining
real FS net \(\phi_i\) on \(B_i\), \(c_1(B_i)\to0\). Let
\(C_i=\phi_i(v_0)\downarrow C=\phi(v_0)\in\mathbb R\). Fix an actual
ample \(A\). Refine the defining indices by rational \(\varepsilon>0\)
for which \[\varepsilon A-B_i\text{ is actual ample}.\tag{10.1}\] For
every \(\varepsilon\), (10.1) holds eventually. The admissible pairs
ordered by increasing \(i\) and decreasing \(\varepsilon\) are directed
and final over the original indices, with \(\varepsilon\to0\): choose a
smaller \(\varepsilon\) first, then an index far enough into its ample
tail. This also works for a stationary original net.

Section multiplication by a generating family of a cleared power of
(10.1) realizes \[a_{i,\varepsilon}=\varepsilon^{-1}(\phi_i-C_i)
\in\operatorname{PSH}(X,A),\qquad a_{i,\varepsilon}(v_0)=0.\tag{10.2}\]
For each divisorial \(v\), continuous evaluation on the compact
normalized actual ample cone has a finite lower bound \(M_v\). Thus
\[\varepsilon M_v\leq\phi_i(v)-C_i\leq0.\tag{10.3}\] Passing to the
refined limit gives \(\phi(v)=C\) on every divisorial point. Section 3.2
puts \(\phi\) into a fixed actual ample \(H\); the constant \(C\) is in
that same cone. G0 divisorial determination proves \(\phi=C\) at all
analytic points. Conversely a constant is an actual trivial-bundle FS
function. This proves \(\operatorname{PSH}(0)=\mathbb R\) directly. It
does not use closure of the sequence \(m\phi\) in a varying or nonample
class.

For the actual trivial bundle on \(n\geq1\) dimensional \(X\), a
constant metric has curvature zero. Its genuine local and global
\(n\)-fold products are zero. Set the reference to zero and define
\(\mathcal E^1_0=\mathbb R\), \(E_0=0\), \(\operatorname{MA}_0=0\).
These are its actual metric products and intersection energy, since
every positive-dimensional factor has class zero. The only positive
target of raw mass zero is \(\mu=0\), and its dual supremum is zero.
Every constant is integrable and solves the equation; any two differ by
a constant. The normalized cone is the singleton zero, comparison is
\(0\leq0\), and the outer energy in every continuous direction is zero.
Envelopes of bounded scalar obstacles are the constants equal to their
infima. No expression divides by \((0^n)\) or converts the zero measure
into a probability. Dimension zero is excluded from this assertion: the
empty product there has unit mass.

This range cannot be enlarged to all zero-volume pseudoeffective classes
by the same uniqueness statement. On \(X=\mathbb P^1\times\mathbb P^1\),
\(L=p_1^*\mathcal O(1)\) is semiample and nontrivial with \(L^2=0\). For
\(b>0\), the actual finite FS potential
\[u_b=\max(\log\|s_0\|,\log\|s_1\|+b)-b\] is normalized and nonconstant.
Along the ray where \(s_1\) vanishes and \(s_0\) is a unit it equals
\(\max(0,b-t)-b\). Its actual metric is pulled back from the first
factor. Smooth FS approximation has two-fold curvature product zero, and
BE uniform continuity gives \(\operatorname{MA}_{\rm raw,L}(u_b)=0\).
Distinct \(b\) yield normalized functions differing by more than
constants. Thus zero measure alone cannot yield uniqueness modulo
constants for a general zero-volume positivity datum. This is a
mathematical obstruction to an unqualified extension, not a solution of
the original axis.

\section{11. Preserved failure, scope and unresolved
statements}\label{preserved-failure-scope-and-unresolved-statements}

The actual108 independent verdict is preserved verbatim in the evidence.
Its counterexample is also proved here. Let \(K\) be the disjoint union
of an isolated \(q\) and a sequence \(p_j\to p_\infty\), all \(p_j\)
isolated. Let \(S=\mathbb N_{>0}\), \(T=\mathbb R\), and
\(\rho(n)=1/n\). Let \(F_n\) consist of extended-continuous functions
with value \(c\) at \(q\), \(-\infty\) at \(p_n\), and \(c-m\)
elsewhere, for \(m\geq1\) and real \(c\). They are stable under
constants. For fixed \(n\), increasing \(m\) gives a decreasing limit
\(g\) equal to zero at \(q\) and \(-\infty\) everywhere else. Thus the
constant outer family \(g_n=g\) with \(\rho(n)\to0\) meets the proposed
varying-family closure hypotheses.

Yet no decreasing net of original family members can have \(\rho\to0\).
After any term from \(F_n\), all later terms must remain in \(F_n\),
because they must keep the value \(-\infty\) at \(p_n\), and every other
family is finite there. Its \(\rho\) is fixed on that entire tail.
Therefore the closure at zero is empty. This disproves BJ v4 Lemma 4.6
as stated. The printed strict shift at a common pole would read
\(-\infty+\epsilon<-\infty+\epsilon'\), which is false; its exponential
difference is zero. Sections 3--4 and 10 do not use that argument or a
hidden varying-class closure behind BJ 4.15, 6.17, 6.24 or 6.31. The
lemma's failure does not disprove Theorem 1. The preserved public
renderer coordinates are BJ v4 L3187--3198@P45 for its statement,
L3222--3241@P45--46 for the defective directedness argument,
L3264--3270@P46 for its closure application, L3550@P49 for the
invocation in 4.15(i), and L4559--4565@P65 for the invocation of 4.15 in
6.17.

The full-big conclusion proved in this paper is an actual equation
theorem for the entire reference-dominated cone for every fixed
\(L=A+D\) datum, with all targets of intrinsic finite energy and raw
mass \((A^n)\). It has no countability, toric, skeletal, support, or
finite-generation restriction. Every construction occurs on the original
smooth projective \(X\). The intrinsic reference is genuinely retained.
Theorem 2 supplies the separately stated trivial-bundle zero-volume
endpoint.

The unchanged original problem is broader if it asks for a class-level
equation over the general pseudoeffective nonample domain or for the
minimal-singularity big cone of mass \(\operatorname{vol}(L)\). Neither
follows from this paper. Concrete remaining mathematical statements are:

\begin{itemize}
\tightlist
\item
  For an arbitrary big rational \(L\), construct and prove the equation
  theory on its minimal-singularity or another generally prescribed
  model-reference cone beyond the divisorial references whose residual
  bundle is actually ample, with its intrinsic mass (in the
  minimal-singularity case, \(\operatorname{vol}(L)\)).
\item
  For arbitrary nonbig pseudoeffective nontrivial rational data,
  identify the correct intrinsic energy, zero or lower-dimensional
  measure datum and an appropriate uniqueness formulation; prove the
  original-carrier operator, compactness, envelopes, differentiation,
  comparison and attainment. Ordinary positive zero-mass measures alone
  cannot retain uniqueness modulo constants, as Section 10 shows.
\item
  If the intended full original theorem includes all reference choices
  or all of \(\operatorname{PSH}(L)\) with one class-intrinsic mass,
  prove compatibility and the necessary replacement inputs for those
  larger cones. No such compatibility has been asserted between distinct
  \((A,D)\) choices here.
\end{itemize}

Accordingly the canonical result is marked \textbf{candidate=false} for
the complete unchanged original P16.5, while retaining Theorems 1--2 and
the direct foundation repairs as substantive new deductions relative to
G0. This conservative scope assessment prevents an entire
prescribed-reference theorem from being advertised as a general
pseudoeffective solution. No further mathematical gap is intentionally
assumed inside the stated prescribed-reference theorem; it remains an
unverified draft with the exact external premises listed below.

\section{Appendix A. External-premise ledger and source
audit}\label{appendix-a.-external-premise-ledger-and-source-audit}

The only task grant is the exact actual-ample G0 body supplied in the
assignment, preserved as \texttt{evidence/block-00.txt}. It is used only
for actual ample rational bundles on the original smooth integral
projective \(X\) and auxiliary actual ample \(H\). Its ample
compactness, envelopes, all-point determination, genuine mixed
intersection operator, energy, weighted unbounded limits, comparison,
intrinsic target criterion, two-sided derivatives, coercivity and
attainment retain every field hypothesis and raw intersection factor.
Its historical bindings identify supplied text, not mathematical
certificates. No other draft, repair proposal, review acceptance, or
general nonample theorem is a grant.

The machine-readable \texttt{external-premise-ledger.json} records exact
saved units, supplied source versions, literal coordinate systems,
supplied bindings, newly saved byte hashes, read scope and recursive
omissions. The substantive premises are these:

\begin{longtable}[]{@{}
  >{\raggedright\arraybackslash}p{(\columnwidth - 4\tabcolsep) * \real{0.3333}}
  >{\raggedright\arraybackslash}p{(\columnwidth - 4\tabcolsep) * \real{0.3333}}
  >{\raggedright\arraybackslash}p{(\columnwidth - 4\tabcolsep) * \real{0.3333}}@{}}
\toprule\noalign{}
\begin{minipage}[b]{\linewidth}\raggedright
ID
\end{minipage} & \begin{minipage}[b]{\linewidth}\raggedright
Precise premise and retained use
\end{minipage} & \begin{minipage}[b]{\linewidth}\raggedright
Primary evidence and limits
\end{minipage} \\
\midrule\noalign{}
\endhead
\bottomrule\noalign{}
\endlastfoot
G0 & The supplied actual-ample theorem in its exact hypotheses; no
extension to nonample classes by deletion & Entire supplied body,
\texttt{evidence/block-00.txt}; 68 supplied dictionary bodies preserved
individually, including their binding keys. Only the relevant ample
statements are used. \\
E0 & Canonical section evaluation and actual FS arithmetic; the
single-net definition & BJ arXiv:1801.08229v4, March 2022; Remark 1.24,
Definitions 2.18 and 4.1, Proposition 2.19 and scaling-ray Lemma 2.20.
Saved pertinent supplied bodies and their literal metadata. No BJ
4.5/4.6 closure is granted. \\
E1 & Openness of the numerical ample cone, finite-dimensional numerical
conventions, actual rational numerical ampleness & BJ conventions
preserved in \texttt{evidence/G0-primary/00.txt}; Fujino--Miyamoto
arXiv:2101.00806v1, 4 January 2021, Theorem 1.4 and its full printed
proof at renderer L47--52 and L179--303. Its complete normal variety and
algebraically closed arbitrary-field hypotheses include our smooth
projective \(X\). Numerical equality preserves the intersection
criterion. Its underlying Birkar augmented-base-locus theorem,
flattening, and classical ampleness references are not recursively
proved here. \\
E2 & Global generation of sufficiently large powers/twists of an actual
ample bundle & Stacks 01Q3, Proposition 28.27.13, complete supplied
mathematical statement/proof; quasi-compact, quasi-separated finite-type
projective \(X\) and coherent ideals satisfy its exact hypotheses. Serre
and Hilbert-polynomial supplied bodies are preserved but are not needed
by the convex-cone construction in Section 2. \\
E3 & Literal section division on a normal variety & BJ v4 Lemma 2.23,
complete pertinent proof in \texttt{evidence/units/primary-01.txt},
original inclusive LF1270--1289 as recorded in its supplied metadata.
The proof uses the standard normal affine height-one regularity
criterion; we use it only on smooth normal original \(X\), with actual
rational bundle and effective rational Cartier divisor. \\
E4 & Smooth regularized FS maxima, continuous actual fundamental-cycle
product and uniform continuity, including trivial valuations & BE
arXiv:1805.01016v3, Sections 7.3 and 8.1: 7.12--7.14, 8.3--8.6 complete
supplied pertinent units. Projective scheme over a complete valued field
and actual semiample rational bundles; the auxiliary bundles used for
global approximation are ample. BE 7.8's nontrivial valuation hypothesis
is respected by not invoking it. The source's printed elementary sign
typo in the integration-by-parts expansion does not replace the correct
difference formula. \\
E5 & Harmonic unit-log frame changes preserve genuine local curvature
products & CLD arXiv:1204.6277v2, continuous metric curvature Section
18.1.2 and local harmonic invariance Section 8.8.5, complete pertinent
supplied units. Used on the boundaryless open original analytification
and its actual isometric line bundles. The local Bedford--Taylor
construction cited inside that work is a published external background;
no fresh recursive proof of its whole theory is claimed. \\
D0 & Failure of the general extended-valued flattening lemma & Supplied
actual108 report, passages and coordinates preserved in
\texttt{evidence/block-02.txt}; Section 11 gives its entire independent
mathematical counterexample. It is a diagnostic, not a positive
premise. \\
\end{longtable}

The actual metadata for each unit take precedence over shortened
descriptions in this table. In particular, coordinates labeled inclusive
one-based LF are on the specified supplied full-text extraction;
coordinates \(Lx@Py\) are public zero-based rendered PDF lines and page
labels. They are not interchanged. The evidence retains the complete
assignment, source-version strings, dictionary keys, supplied source
paths as documentary strings, and original body and source hashes. No
path pointing to another worker was opened; only the inline data and
newly public primary responses were used. No unavailable source PDF
checksum is invented.

The printed BJ nonample homogeneous and Siu units, the BJ+
arXiv:2206.07183v2 envelope proof and its dependence on the disputed
closure chain, the unaccepted105 body and register, and continuations
109--114 are preserved as supplied data. None supplies an unproved
nonample conclusion to this paper. Sections 3--4 reconstruct the
proposed foundation repair rather than cite continuation114. The
fixed-ample flattening explanation does not treat continuation113 as a
grant. The negligible-set proof involving Choquet extraction is not
used; Section 6.2 gives the needed polar tests with its capacity weight
and sufficient exponent explicit.

\section{Appendix B. Reproduction, checks and evidence
limits}\label{appendix-b.-reproduction-checks-and-evidence-limits}

\texttt{input.json} is kept byte-for-byte.
\texttt{assignment-original.txt} is the exact decoded assignment string
and \texttt{task-original.json} preserves the exact task value.
\texttt{evidence/block-index.json} records all extracted supplied
blocks; their complete byte hashes and original metadata remain in the
evidence manifest. The extraction never dereferences a supplied path.
\texttt{rerun.py} regenerates the blocks and individual units, verifies
retained bindings where their byte delimitation agrees, reports
mismatches instead of rebranding them as proofs, regenerates the premise
ledger, checks the exact five-key status schema and backward note
supports, and rebuilds this Markdown paper into TeX and PDF with the
installed Pandoc and pdfLaTeX.

Run \texttt{python3\ rerun.py} in this workspace. \texttt{checks.py}
performs only finite rational diagnostic checks of the polytope formula,
the rounding inequalities, a finite shadow of the failure example, and
the zero-volume ray calculation. Its domains and counts are recorded in
\texttt{checks-output.json}; they do not prove the net, analytic,
algebraic, or variational statements. The proof is the arguments in the
paper relative to the stated premises. Commands, argv, outputs,
executable versions, platform and hashes are saved locally in
\texttt{run-log.json}, \texttt{environment.json}, \texttt{build/}, and
\texttt{artifact-hashes.json}. No tools are installed and no heavy local
computation is launched; no operator supervised compute facility was
provided or needed. Compilation and finite checks establish only
reproducibility and diagnostics.

\end{document}
